\documentclass[10pt]{article}
\usepackage[T1]{fontenc}
\usepackage[utf8]{inputenc}
\usepackage[french]{babel}
\usepackage{amsmath,amssymb,mathrsfs,mathpazo}
\usepackage{geometry}
\geometry{a4paper,margin=2cm}
\usepackage{graphicx,array,tabularx,tabularray,colortbl,multirow,multicol,booktabs,enumitem,pifont,ulem,fancybox,diagbox,cancel}
\usepackage{pstricks,pstricks-add,pst-plot,pst-tree,pst-node,pst-text,pst-func,pst-math,pst-bspline,pst-3dplot,pst-eucl,pst-solides3d}
\let\pscalculate\undefined
\usepackage{pst-bezier}
\usepackage{tikz,pgfplots,tkz-tab}
\usetikzlibrary{arrows,arrows.meta,patterns,calc,positioning,trees,decorations.markings,shapes.geometric,shapes.misc,shapes.arrows,intersections,decorations.pathmorphing,decorations.pathreplacing,backgrounds,shadows,babel}
\pgfplotsset{compat=1.18}
\usepgfplotslibrary{fillbetween}
\usepackage[np]{numprint}
\usepackage{listings,makecell,pas-tableur,twoopt}
\let\barre\undefined
\usepackage[pstricks]{tablvar}
\definecolor{darkgreen}{RGB}{0,100,0}
\providecommand{\red}{\color{red}}
\providecommand{\blue}{\color{blue}}
\providecommand{\green}{\color{green}}
\providecommand{\black}{\color{black}}
\providecommand{\white}{\color{white}}
\providecommand{\gray}{\color{gray}}
\providecommand{\cyan}{\color{cyan}}
\providecommand{\magenta}{\color{magenta}}
\providecommand{\yellow}{\color{yellow}}
\providecommand{\violet}{\color{violet}}
\providecommand{\orange}{\color{orange}}
\providecommand{\euro}{\texteuro}
\providecommand{\EUR}{\texteuro}
\providecommand{\ieme}{\textsuperscript{e}}
\providecommand{\degres}{\ensuremath{^\circ}}
\providecommand{\degree}{\ensuremath{^\circ}}
\providecommand{\vv}[1]{\overrightarrow{#1}}
\providecommand{\Coord}[1]{#1}
\providecommand{\decosix}{}
\providecommand{\starredbullet}{\ensuremath{\bullet}}
\pagestyle{empty}
\setlength{\parindent}{0pt}
\usepackage{adjustbox}\DeclareUnicodeCharacter{2236}{:}\DeclareUnicodeCharacter{2190}{\ensuremath{\leftarrow}}
\providecommand{\R}{\mathbb{R}}
\providecommand{\N}{\mathbb{N}}
\providecommand{\Z}{\mathbb{Z}}
\providecommand{\Q}{\mathbb{Q}}
\providecommand{\D}{\mathbb{D}}
\providecommand{\C}{\mathbb{C}}
\providecommand{\e}{{\mathrm{e}}}
\providecommand{\vect}[1]{\overrightarrow{#1}}
\providecommand{\vectt}[1]{\overrightarrow{\mathrm{#1}}}
\providecommand{\vvt}[1]{\overrightarrow{\mathrm{#1}}}
\providecommand{\np}[2][]{{#2~#1}}
\providecommand{\vv}[1]{{\overrightarrow{#1}}}
\renewcommand{\d}{{\mathrm{d}}}
\providecommand{\degres}{{^\circ}}
\providecommand{\degre}{{^\circ}}
\providecommand{\blue}{\color{blue}}
\providecommand{\red}{\color{red}}
\providecommand{\enskip}{\quad}
\providecommand{\llbracket}{\lbrack\!\lbrack}
\providecommand{\rrbracket}{\rbrack\!\rbrack}
\providecommand{\euro}{{\text{€}}}
\providecommand{\no}{{\text{n°}}}
\providecommand{\interval}[2]{{[#1;#2]}}
\providecommand{\intervFO}[2]{{[#1;#2[}}
\providecommand{\dotfill}{\ldots}

\providecommand{\ds}{\displaystyle}
\providecommand{\textding}[1]{#1}
\providecommand{\Oij}{\(\mathrm{O};\vec{\imath},\vec{\jmath}\)}
\providecommand{\Oijk}{\(\mathrm{O};\vec{\imath},\vec{\jmath},\vec{k}\)}
\providecommand{\Ouv}{\(\mathrm{O};\vec{u},\vec{v}\)}
\providecommand{\intervFF}[2]{\ensuremath{[#1;#2]}}
\renewcommand{\euro}{\ifmmode\text{\texteuro}\else\texteuro\fi}
\DeclareUnicodeCharacter{00B0}{\ifmmode^{\circ}\else\textdegree\fi}
\DeclareUnicodeCharacter{20AC}{\ifmmode\text{\texteuro}\else\texteuro\fi}
% Énoncés : Éducation nationale. Transcription et corrigés : APMEP, Denis Vergès et contributeurs. Découpage et classement : CoopMaths, Sébastien Lozano, Éric Elter et contributeurs. Adaptation de la présentation : Maths Explorer.
% CC BY-SA : https://coopmaths.fr/www/about/
% https://coopmaths.fr/alea/static/bac/2023/tex/bac_2023_03_sujet2_metropole_3_cor.tex
\begin{document}
\noindent\textit{APMEP / CoopMaths. CC BY-SA. Présentation adaptée par Maths Explorer.}\par\medskip
\medskip

Dans l'espace rapporté à un repère orthonormé \Oijk, on considère :

\begin{itemize}
\item[$\bullet~~$]le plan $\mathcal{P}_1$ dont une équation cartésienne est $2x + y - z +2 = 0$,
\item[$\bullet~~$]le plan $\mathcal{P}_2$ passant par le point B(1~;~1~;~2) et dont un vecteur normal est $\vect{n_2}\begin{pmatrix}1\\-1\\1\end{pmatrix}$.
\end{itemize}



\begin{enumerate}
\item 
	\begin{enumerate}
		\item  Le vecteur $\vect{n_1}\;\begin{pmatrix} 2 \\1 \\ -1\end{pmatrix}$ est normal au plan $\mathcal{P}_1$ d'équation $2x+y-z+2=0$.
		
		\item 
		


$\vect{n_1}.\vect{n_2} = 2\times 1 + 1\times (-1) +(-1)\times 1 = 0$ donc $\vect{n_1}$ et $\vect{n_2}$ sont orthogonaux. 

On en déduit que les plans $\mathcal{P}_1$ et $\mathcal{P}_2$ sont perpendiculaires. 
	\end{enumerate}
\item 
	\begin{enumerate}
		\item 
Le plan $\mathcal{P}_2$ a pour vecteur normal $\vect{n_2}\begin{pmatrix}1\\-1\\1\end{pmatrix}$, donc il a une équation cartésienne de la forme $x-y+z+d=0$ où $d$ est un réel à déterminer.

$\text{B} \in \mathcal{P}_2$ donc $x_{\text B}-y_{\text B}+z_{\text B}+d=0$, autrement dit $1 - 1 +2 +d=0$ donc $d=-2$.

Le plan $\mathcal{P}_2$ a donc pour équation cartésienne $x-y+z-2=0$.

		\item On note $\Delta$ la droite dont une représentation paramétrique est : $\left\{\begin{array}{l !{=} l}
x&0\\y & -2 + t\\ z&\phantom{2+t}t
\end{array}\right.,\quad t \in  \R$.

On cherche l'intersection des plans $\mathcal{P}_1$ et $\mathcal{P}_2$ est l'ensemble des points de coordonnées $(x\;;\;y\;;\;z)$ vérifiant le système:
$\left \{
\begin{array}{l !{=} l}
2x+y-z +2 & 0\\
x - y + z - 2&0
\end{array}
\right .$.

$\left \{
\begin{array}{l !{=} l}
2x+y-z +2 & 0\\
x - y + z - 2&0
\end{array}
\right .
\iff
\left \{
\begin{array}{l !{=} l}
2x+y & z-2\\
x - y &-z+2
\end{array}
\right .
\iff
\left \{
\begin{array}{l !{=} l}
3x&0\\
2x+y & z-2
\end{array}
\right .$

On aboutit donc à $x=0$, $y=z-2$ et $z$ quelconque égal à $t$.

Les plans $\mathcal{P}_1$ et $\mathcal{P}_2$ ont donc pour intersection la droite de représentation paramétrique 
$\left \{
\begin{array}{l !{=} r l}
x & 0\phantom{+t}\\
y & -2+t & t\in\R\\
z & t
\end{array}
\right.$, c'est-à-dire la droite $\Delta$.
	\end{enumerate}
	\end{enumerate}
	
On considère le point A(1~;~1~;~1) et on admet que le point A n'appartient ni à $\mathcal{P}_1$ ni à $\mathcal{P}_2$.

On note H le projeté orthogonal du point A sur la droite $\Delta$.

\begin{enumerate}[resume]
\item On rappelle que, d'après la question 2. b, la droite $\Delta$ est l'ensemble des points M$_t$ de coordonnées $(0~;~-2 + t~;~t)$, où $t$ désigne un nombre réel quelconque.

	\begin{enumerate}
		\item 
		
$\text{AM}_t^2 = (0-1)^2 + (-2+t-1)^2 +(t-1)^2 = 1 + \left (9 -6t +t^2\right ) +\left (t^2 -2t +1\right ) = 2t^2 -8t +11$

Donc $\text{AM}_t = \ds\sqrt{2t^2-8t+11}$.
		
		\item 
Le point H est le projeté orthogonal de A sur la droite $\Delta$, donc la longueur AH  réalise le minimum des longueurs AM$_t$ où M$_t$ est un point de $\Delta$.

Il faut donc chercher le minimum de $ \ds\sqrt{2t^2-8t+11}$, donc le minimum de $2t^2-8t+11$.

D'après les propriétés de la fonction du second degré, le minimum de $f(x)=ax^2+bx+c$ quand $a>0$, est réalisé pour $x=-\dfrac{b}{2a}$ et vaut $f\left (-\dfrac{b}{2a}\right )$.

Donc le minimum de $2t^2 - 8t + 11$ est réalisé pour $t=-\dfrac{-8}{2\times 2}=2$, et  vaut $2\times 2^2 -8\times 2 + 11 = 3$.

On en déduit que $\text{AH}=\ds\sqrt{3}$.
		
		
	\end{enumerate}
\item On note $\mathcal{D}_1$ la droite orthogonale au plan $\mathcal{P}_1$ passant par le point A et H$_1$ le projeté orthogonal du point A sur le plan $\mathcal{P}_1$.
	\begin{enumerate}
		\item Déterminer une représentation paramétrique de la droite $\mathcal{D}_1$.
		
La droite $\mathcal{D}_1$ est orthogonale au plan $\mathcal{P}_1$ donc le vecteur $\vect{n_1}$, normal au plan $\mathcal{P}_1$ est un vecteur directeur de la droite $\mathcal{D}_1$. De plus la droite $\mathcal{D}_1$ passe par le point A. Elle a donc pour représentation paramétrique:

$\left \{
\begin{array}{l !{=} r l}
x & x_{\text A} + 2t\\
y & y_{\text A} +t & t\in \R\\
z & z_{\text A} -t
\end{array}
\right .$
c'est-à-dire
$\left \{
\begin{array}{l !{=} r l}
x & 1 + 2t\\
y & 1 +t & t\in \R\\
z & 1 -t
\end{array}
\right .$		
		
		\item 
La droite $\mathcal{D}_1$ est orthogonale au plan $\mathcal{P}_1$ donc le projeté orthogonal de A sur $\mathcal{P}_1$ est le point d'intersection de $\mathcal{D}_1$ et de $\mathcal{P}_1$; ses coordonnées vérifient le système:

$\left \{
\begin{array}{r !{=} l}
x & 1 + 2t\\
y & 1 +t \\
z & 1 -t\\
2x +y-z+2 & 0
\end{array}
\right .$

Donc $t$ vérifie $2(1+2t) +(1+t) -(1-t) +2=0$, soit $2+4t+1+t-1+t+2=0$, ce qui donne $t=-\dfrac{2}{3}$.

$x=1+2t = 1-\dfrac{4}{3} =-\dfrac{1}{3}$, $y=1+t = 1-\dfrac{2}{3}=\dfrac{1}{3}$ et $z=1-t= 1+\dfrac{2}{3}=\dfrac{5}{3}$

Le point H$_1$ a donc pour coordonnées $\left (-\dfrac{1}{3}\;;\; \dfrac{1}{3}\;;\;\dfrac{5}{3} \right )$.
		
	\end{enumerate}
\item Soit H$_2$ le projeté orthogonal de A sur le plan $\mathcal{P}_2$.

\begin{minipage}{0.48\linewidth}

On admet que H$_2$ a pour coordonnées $\left(\dfrac43~;~\dfrac23~;~\dfrac43\right)$
et que H a pour coordonnées (0~;~0~;~2). 

Sur le schéma ci-contre, les plans $\mathcal{P}_1$ et $\mathcal{P}_2$
sont représentés, ainsi que les points A, H$_1$, H$_2$, H. 
\end{minipage}\hfill
\begin{minipage}{0.48\linewidth}
\scalebox{0.8}{
\psset{unit=1cm}
\begin{pspicture}(7.4,8.6)

\pspolygon(0,2.7)(5.5,0.5)(7.2,4.3)(1.7,6.5)
\psline(2.5,0)(2.5,4.7)(4.2,8.4)(4.2,5.3)
\psline(1.8,0)(3.1,2.7)
\psline[linestyle=dashed](3.1,2.7)(3.6,3.62)
\psline[linestyle=dashed](4.2,5.3)(4.2,1.1)
\psline(3.6,3.62)(5.6,8)
\psline(4.2,1.1)(3.7,0)
\pspolygon(3.1,2.7)(4.6,2.3)(4.6,3.4)(3.1,3.8)
\uput[l](2,0.4){$\Delta$}\rput(0.4,2.8){$\mathcal{P}_2$}\rput(4,7.4){$\mathcal{P}_1$}
\uput[dl](3.1,2.7){H}\uput[dr](4.6,2.3){H$_2$}\uput[ur](4.6,3.4){A}\uput[ul](3.1,3.8){H$_1$}
\end{pspicture}
}
\end{minipage}



Le vecteur $\vect{\text{AH}_1}$ a pour coordonnées
$\begin{pmatrix}
-\frac{1}{3}-1 \\ \frac{1}{3}-1 \\ \frac{5}{3}-1
\end{pmatrix}$
=
$\begin{pmatrix}
-\frac{4}{3} \\ -\frac{2}{3} \\ \frac{2}{3}
\end{pmatrix}$.

Le vecteur $\vect{\text{H}_2\text{H}}$ a pour coordonnées
$\begin{pmatrix}
0- \frac{4}{3}\\ 0- \frac{2}{3} \\ 2- \frac{4}{3}
\end{pmatrix}$
=
$\begin{pmatrix}
-\frac{4}{3} \\ -\frac{2}{3} \\ \frac{2}{3}
\end{pmatrix}$.

$\vect{\text{AH}_1} = \vect{\text{H}_2\text{H}}$ donc la quadrilatère AH$_1$HH$_2$ est un parallélogramme.

La droite (AH$_1$) est orthogonale au plan $\mathcal{P}_1$ et H$_1$  appartient à ce plan; donc (AH$_1$) est perpendiculaire à toutes les droites de $\mathcal{P}_1$ passant par H$_1$, en particulier la droite (HH$_1$).

Le parallélogramme AH$_1$HH$_2$  possède un angle droit donc c'est un rectangle.

\end{enumerate}

\bigskip

\end{document}
